A Heegaard splitting of a closed 3-manifold can be represented by a collection of disjoint simple closed curves on the boundary of a handlebody.  In the previous section, we examined the girth of simple closed curves on the boundary of a handlebody.  In this section we turn to girth for Heegaard splittings.  We define the algebraic and geometric girths of a Heegaard splitting.  We present Conjecture~\ref{ag girth conjecture} and prove that it is equivalent to the Girth Conjecture.  We also present some connections to the study of 3-manifolds with finite fundamental group.

Let $M$ be an orientable closed 3-manifold.  A \emph{Heegaard splitting} 
$H\cup H^\prime$ of $M$ is an embedding of a handlebody $H$ into $M$ so that $M-H$ is another handlebody $H^\prime$.  The intersection of the handlebodies in $M$ is the surface $\bd H$.  The order matters in our definition, so the splitting $H\cup H^\prime$ is different from $H^\prime\cup H$.  

Let $H\cup H^\prime$ be a Heegaard splitting of $M$.  The fundamental 
group $\pi_1(M)$ equals the quotient of $\pi_1(H)$ by some normal 
subgroup of relators $\rr$.
Each element of $\rr$ can be represented by an immersed 
curve in $\partial H$ that bounds a disc in $H^\prime$.
By the following lemma, the trivial element of $\rr$ can 
always be represented by an essential closed curve in $\bd H$.
This is why we require primitive curves $\gamma$ to be nontrivial in our 
statement of the Girth Conjecture.
The following lemma appears as Theorem~10.2 in \cite{P}.
We present a short proof here for clarity.

\begin{lemma}
\label{trivial relator}
Let $M$ be a closed orientable $3$-manifold.  Let $H\cup H^\prime$ 
be a Heegaard splitting of $M$ with genus strictly bigger than one.  
Then there is an essential closed curve in $\partial H$ that 
contracts in both $H$ and $H^\prime$.
\end{lemma}

\begin{proof}
Pick a basepoint $x$ in $\partial H$.  Choose an essential loop
$g$ (resp. $h$) based at $x$ that contracts in $H$ (resp. $H^\prime$).
Since $\genus(\partial H)>1$, we can choose $g$ and $h$ not to commute
in $\partial H$.  The commutator $[g,h]$ is an essential 
loop in $\partial H$ that contracts in both $H$ and $H^\prime$.
\end{proof}

Define the \emph{algebraic girth} of $H\cup H^\prime$ to be 
the minimum girth in $H$ over all nontrivial elements of $\rr$.  
When $\rr$ is trivial, we define the algebraic girth to be zero.  
Observe that $\rr$ is trivial if and only if $M$ is a connect 
sum of $g$ $S^2 \times S^1$'s, where $g$ is the genus of $H$.
The definition of algebraic girth is not symmetric because
it depends on a choice of one of the two handlebodies.  
It is not clear \emph{a priori} whether the algebraic girths of $H\cup H^{\prime}$ and $H^\prime \cup H$ are the same .

By the following lemma, higher genus Heegaard splittings of 3-manifolds with finite fundamental group have algebraic girth zero.

\begin{lemma}
\label{finite implies girth zero}
Let $M$ be a closed orientable 3-manifold.  Let $H\cup H^{\prime}$ be a Heegaard splitting of $M$.  If $M$ has finite fundamental group and $\genus(H)\geq 2$, then the algebraic girth of $H\cup H^{\prime}$ is zero.
\end{lemma}

\begin{proof}
Let $\dd$ be a disc system for $H$.  Let $\{ a, b, \ldots \}$ be the corresponding generating system for $\pi_{1}(H)$.  Consider the generator $a$.  Represent $a$ as a closed curve disjoint from the disc corresponding to the generator $b$.  Since $\pi_1(M)=\pi_1(H)/\rr$ is finite, there is an $n\geq 1$ so that $a^n$ is trivial in $\pi_1(M)$.  Since $a^n$ has girth zero in $H$, the splitting $H\cup H^\prime$ has algebraic girth zero.
\end{proof}


Let $\scc$ denote the collection of simple closed curves in $\partial H$ that bound a disc in $H^\prime$.  Define the \emph{geometric girth} of $H\cup H^\prime$ to be the minimum girth in $H$ over all essential elements of $\scc$.  By Lemma~\ref{one embedded curve}, the geometric girth equals the minimum intersection number of a pair of essential discs on either side of $\partial H$.  For example, a splitting is weakly reducible if and only if it has geometric girth zero.  Observe that geometric girth is symmetric with respect to $H$ and $H^{\prime}$.

The algebraic girth of a Heegaard splitting $H\cup H^\prime$ is less than or equal to its geometric girth.  We state this as a corollary of the following lemma.

\begin{lemma}
\label{girth zero}
Let $M$ be a closed orientable $3$-manifold.  Let $H\cup H^\prime$ 
be a Heegaard splitting of $M$.  If $H\cup H^\prime$ is reducible, then the algebraic girth of $H\cup H^\prime$ is zero.
\end{lemma}

Recall that a Heegaard splitting $H\cup H^{\prime}$ is reducible if there exists an essential  simple closed curve $\gamma$ in $\bd H$ that bounds a disc in both $H$ and $H^{\prime}$.

\begin{proof}
Since $H\cup H^{\prime}$ is reducible, choose a simple closed curve $\gamma$ in $\bd H$ that bounds a disc $D$ in $H$ and a disc $D^\prime$ in $H^\prime$.  Complete the disc $D^\prime$ in $H^\prime$ to a disc system $\dd^{\prime}$ in $H^\prime$.  Then the curves in $\bd \dd^{\prime}$ normally generate the group of relators $\rr$.  If all the curves in $\bd\dd^{\prime}$ are trivial in $\pi_1(H)$, then the group of relators $\rr$ is trivial; and the algebraic girth is zero by definition.

Otherwise, at least one curve in $\bd\dd^{\prime}$ is nontrivial in $\pi_1(H)$.  Consider such a curve.  This curve is disjoint from $D$, so it has girth zero in $H$.  It represents a nontrivial element of $\rr$, so the splitting has algebraic girth zero.
\end{proof}

We now prove a corollary of Lemma~\ref{girth zero}.

\begin{corollary}
\label{easy inequality}
Let $M$ be a closed orientable 3-manifold.  Let $H\cup H^\prime$ be a Heegaard splitting of $M$.  The algebraic girth of $H\cup H^\prime$ is less than or equal to its geometric girth.
\end{corollary}

\begin{proof}
Choose a simple closed curve $\gamma$ in $\bd H$ that realizes 
the geometric girth of $H\cup H^\prime$. Then $\gamma$ bounds a disc $D^\prime$ in $H^\prime$, and the girth of $\gamma$ in $H$ equals the geometric girth of $H\cup H^\prime$.

If $\gamma$ does not bound a disc in $H$, then $\gamma$ represents a nontrivial element of $\rr$.  Viewing $\gamma$ as an element of $\rr$, we conclude that the algebraic girth of $H\cup H^\prime$ is less than or
equal to the geometric girth.  If $\gamma$ does bound a disc $D$ in $H$, then $H\cup H^{\prime}$ is reducible.  Applying Lemma~\ref{girth zero}, we are done.
\end{proof}

The Girth Conjecture is equivalent to the assertion that the
reverse inequality is also true, except for a special case.  
We state the assertion here as a conjecture, and we prove its
equivalence below as Proposition~\ref{a and b}.

\begin{conjecture}
\label{ag girth conjecture}
Let $H\cup H^\prime$ be a Heegaard splitting of a closed orientable 
$3$-manifold.  If $H\cup H^\prime$ is not a genus two splitting of a genus two spherical manifold, then the algebraic and geometric girths of $H\cup H^\prime$ are equal.
\end{conjecture}

This is not necessary for our discussion, but the algebraic and geometric girths of genus two splittings of genus two spherical manifolds are not equal.  These splittings have algebraic girth zero and geometric girth two.  They have algebraic girth zero by Lemma~\ref{finite implies girth zero}.  They have geometric girth two for the following reasons.

Splittings of Seifert-fibered spaces have been classified in~\cite{MS}.  These splittings can be verified to have geometric girth at least two or less.  Genus two Heegaard splittings of geometric girth zero are all connect sums of Lens spaces, so their girth cannot be zero.  Finally, geometric girth one splittings of genus two or more do not exist by Lemmas~\ref{one embedded curve} and~\ref{geometric girth one}.

We prove Conjecture~\ref{ag girth conjecture} is equivalent to the Girth Conjecture at the end of this section.  For reference, we state their equivalence as a proposition here.

\begin{proposition}
\label{a and b}
Let $H\cup H^{\prime}$ be a Heegaard splitting of a closed orientable 3-manifold.  The Girth Conjecture is true for $H\cup H^{\prime}$ if and only if Conjecture~\ref{ag girth conjecture} is true for $H\cup H^{\prime}$.
\end{proposition}

It is often easier to use Conjecture~\ref{ag girth conjecture} in proofs than to cite the Girth Conjecture directly.  Here are three examples.

\begin{lemma}
\label{reducible}
Let $M$ be a closed orientable 3-manifold.  Let $H\cup H^\prime$ be a Heegaard splitting of $M$.  If $M$ is reducible, or more generally if $H\cup H^{\prime}$ is reducible, then the Girth Conjecture is true for $H\cup H^{\prime}$.
\end{lemma}

\begin{proof}
By~\cite{Ha}, the splitting $H\cup H^{\prime}$ is reducible if $M$ is reducible.  So consider a reducible Heegaard splitting $H\cup H^{\prime}$.

Since $H\cup H^{\prime}$ is reducible, it has geometric girth zero.  By Lemma~\ref{girth zero}, the splitting has algebraic girth zero.  Since the algebraic and geometric girths are equal, Conjecture~\ref{ag girth conjecture} is true for $H\cup H^{\prime}$.  By Proposition~\ref{a and b}, the Girth Conjecture is true for $H\cup H^{\prime}$.
\end{proof}

The Girth Conjecture for 3-manifolds with finite fundamental group is equivalent to the spherical space form conjecture.  The Geometrization Conjecture, then, implies the Girth Conjecture for Heegard splittings of 3-manifolds with finite fundamental group.  Conversely, the Girth Conjecture would provide another proof of the Geometrization Conjecture for 3-manifolds with finite fundamental group.

\begin{proposition}
\label{finite fundamental group}
Let $M$ be a closed orientable 3-manifold with finite fundamental group.  Let $H\cup H^\prime$ be an irreducible Heegaard splitting of $M$.  The 3-manifold $M$ is spherical if and only if the Girth Conjecture is true for $H\cup H^\prime$.
\end{proposition}

\begin{proof}
\textbf{S $\Rightarrow$ GC:}  
Let $M$ be a spherical 3-manifold.  By Lemma~\ref{reducible}, if $H\cup H^{\prime}$ is reducible, then the Girth Conjecture is true for $H\cup H^{\prime}$.  So assume the splitting $H\cup H^{\prime}$ is irreducible.

Irreducible Heegaard splittings of spherical 3-manifolds have been classified, and they are all genus two or less.  If $H$ has genus two, then $H\cup H^{\prime}$ satisfies Conjecture~\ref{ag girth conjecture}.  By Proposition~\ref{a and b}, it must also satisfy the Girth Conjecture.

So assume $H$ has genus one.  In this case, $H\cup H^{\prime}$ is a standard genus one splitting of a Lens space.  These have equal algebraic and geometric girths.  Consequently, they satisfy Conjecture~\ref{ag girth conjecture}, and hence the Girth Conjecture.

\textbf{GC $\Rightarrow$ S:}  
Let $H\cup H^{\prime}$ be a Heegaard splitting of $M$ that satisfies the Girth Conjecture.  By Proposition~\ref{a and b}, Conjecture~\ref{ag girth conjecture} is also true for $H\cup H^{\prime}$.

If the genus of $H$ is one, then $M$ is a Lens space and so is a spherical 3-manifold.  So assume that the genus of $H$ is at least two.  By Conjecture~\ref{ag girth conjecture}, either $M$ is a genus two spherical 3-manifold, or the algebraic and geometric girths of $H\cup H^\prime$ are equal.  In the former case, we are done.  So assume that the algebraic and geometric girths are equal.

By Lemma~\ref{finite implies girth zero}, the splitting $H\cup H^{\prime}$ has algebraic girth zero.  Therefore $H\cup H^{\prime}$ has geometric girth zero.  In other words, $H\cup H^\prime$ is weakly reducible.  By a theorem of Casson and Gordon \cite{CG}, the 3-manifold $M$ is either Haken or the splitting is 
reducible.  Since $M$ has finite fundamental group, the manifold cannot 
be Haken.  Thus $H\cup H^\prime$ is reducible.
\end{proof}

Recall Lemma~\ref{geometric girth one} which states that no simple closed curve on the boundary of a higher genus handlebody has girth one.  Equivalently, any Heegaard splitting of geometric girth one is the standard genus one splitting of the 3-sphere.  The same is true for algebraic girth.  Together these facts imply that the Girth Conjecture is true for all Heegaard splittings of algebraic girth one.  We state this as a proposition below.  For comparison, note that Conjecture~\ref{ag girth conjecture} and the Girth Conjecture are not known even for Heegaard splittings of algebraic girth zero.

\begin{proposition}
\label{algebraic girth one}
Let $H\cup H^\prime$ be a Heegaard splitting of a closed 3-manifold.  If $H\cup H^\prime$ has algebraic girth one, then the splitting is the standard genus one splitting of the 3-sphere.  In particular, the Girth Conjecture holds for $H\cup H^\prime$.
\end{proposition}

Note that Proposition~\ref{algebraic girth one} also follows from Lemma~\ref{group girth one}.

\begin{proof}
Let $H\cup H^\prime$ be a Heegaard splitting of a closed orientable 3-manifold.  Let $\rr$ denote the group of relators, so that $\pi_1(M)=\pi_1(H)/\rr$.  Assume that $\rr$ has a relator of girth one.

It suffices to show that either $H\cup H^\prime$ is a genus one splitting of the 3-sphere, or $H\cup H^\prime$ has algebraic girth zero.  For if $H\cup H^{\prime}$ were the genus one splitting of the 3-sphere, then it would have geometric girth one.  This would imply Conjecture~\ref{ag girth conjecture} for $H\cup H^{\prime}$.  By Proposition~\ref{a and b}, this would also imply the Girth Conjecture for $H\cup H^{\prime}$.

Let $r$ be an element of $\rr$ of girth one.  Let $\kappa$ be a curve in $H$ representing $r$.  Then $\kappa$ can be made to intersect some essential disc $D$ in $H$ exactly once.  The disc $D$ must be non-separating.

Let $n$ be the genus of $H$.  If $n$ equals one, then $\pi_1(M)$ is trivial and $H\cup H^\prime$ represents the 3-sphere.  So assume that $n$ is greater than one.  By Lemma~\ref{genus two}, the genus of $H$ cannot equal two, so we have $n\geq 3$.

Extend $D$ to a disc system $\dd$ for $H$.  Let $\{a_{1}, a_{2}, \ldots a_{n}\}$ denote the corresponding generating system for $\pi_{1}(H)$, with $a_{1}$ corresponding to $D$.  Homotope $\kappa$ to minimize intersection with $\dd$.  By Corollary~\ref{subcollection}, the curve $\kappa$ intersects $D$ exactly once.  Let $r=a_{1}\cdot w(a_{2}, a_{3}, \ldots a_{n})$ denote the corresponding word in the generators, where $a_{1}$ does not occur in the word $w$.  Then $r$ is cyclically reduced.

Choose a disc system $\dd^{\prime}$ for $H^\prime$.  Then $\bd\dd^{\prime}$ gives a collection of $n$ relators $g_i$ that normally generate $\rr$.
Multiplying by conjugates of $r$, we can eliminate every occurrence of $a_{1}$ from the $g_{i}$.  This gives $n$ new relators $g_i$ that, together with $r$,
normally generate $\rr$.  If one of the $g_{i}$, say $g_{1}$, is nontrivial in $\pi_{1}(H)$, then $H\cup H^{\prime}$ has algebraic girth zero.  This follows because $g_{1}$ would be a nontrivial relator represented by a curve disjoint from $D$.

Assume that all of the $g_{i}$ are trivial in $\pi_{1}(H)$.  Then $\rr$ is normally generated by the relator $r$.  We prove that $\pi_{1}(M)$ equals the free group generated by $\{ a_{2}, a_{3}, \ldots a_{n}\}$.  Let $\Phi$ be the natural map $\Phi:\langle a_{2}, a_{3}, \ldots a_{n} \rangle\rightarrow\pi_{1}(H)/\rr=\langle a_{1}, a_{2}, \ldots a_{n} \rangle/ (r)$.  The map $\Phi$ is surjective because $a_{1}=w(a_{2}, a_{3}, \ldots a_{n})^{-1}$ in $\pi_{1}(H)/\rr$.  By the Freiheitssatz~\cite[Thm. 4.10]{MKS}, every nontrivial relator in the normal subgroup generated by $r$ must involve $a_{1}$.  It follows that $\Phi$ is injective.

Consequently, $\pi_{1}(M)$ is isomorphic to a free group on $n-1$ letters.  Since $n\geq 3$, the fundamental group $\pi_{1}(M)$ is equal to a nontrivial free product $A\star B$.  By Kneser's Conjecture~\cite[Theorem II.A.3]{St2}, the 3-manifold $M$ is reducible.  By~\cite{Ha}, this implies the Heegaard splitting $H\cup H^{\prime}$ is reducible.  Finally, applying Lemma~\ref{girth zero}, the algebraic girth of $H\cup H^{\prime}$ is zero.
\end{proof}

We now prove Proposition~\ref{a and b}.

\begin{proof}[Proof of Proposition~\ref{a and b}]
Let $H\cup H^\prime$ be a Heegaard splitting of a closed
orientable 3-manifold $M$.

\textbf{\ref{ag girth conjecture} $\Rightarrow$ GC:}  
Assume Conjecture~\ref{ag girth conjecture} for $H\cup H^{\prime}$.
Let $\gamma$ be a closed curve in $\bd H$ that is primitive 
in $\pi_1(H)$ and contracts in $H^\prime$.  It suffices to construct a simple closed curve $\gamma^{\prime}$ in $\bd H$ that bounds a disc in $H^{\prime}$ and satisfies $\girth(\gamma^{\prime}) \leq \girth(\gamma)$.

Since $\gamma$ is primitive, the curve $\gamma$ represents a nontrivial element of $\rr$.  In particular, the algebraic girth of $H\cup H^\prime$
is less than or equal to $\girth(\gamma)$.
Choose a simple closed curve $\gamma^\prime$ in $\bd H$ that 
realizes the geometric girth of $H\cup H^\prime$.  
We will show that $\girth(\gamma^{\prime}) \leq \girth(\gamma)$.  By Conjecture~\ref{ag girth conjecture}, either $H\cup H^{\prime}$ is a genus two splitting of a genus two spherical manifold, or the algebraic and geometric girths are equal.

Assume first that the algebraic and geometric girths are equal.  Then $\girth(\gamma^{\prime})$ equals the algebraic girth of $H\cup H^\prime$.  Because the algebraic girth of $H\cup H^\prime$ is less than or equal to $\girth(\gamma)$, we are done in this case.

Assume next that $H\cup H^\prime$ is a genus two splitting of 
a genus two spherical manifold.  Then $\pi_1(M)$ is not a 
cyclic group.  Moreover, the geometric girth of $H\cup H^\prime$
is exactly two.  This is by our remarks following the statement of Conjecture~\ref{ag girth conjecture}.  Hence $\girth(\gamma^\prime)=2$.  It suffices to show that $\gamma$ cannot have girth zero or one.

In a genus two handlebody, the only primitive curves of girth zero 
correspond to group elements of the form $a$, where $\langle a, b \rangle$ is a 
generating system for $\pi_1(H)$.  The curve $\gamma$ cannot be of this form because this would imply $\pi_1(M)$ is cyclic.  Thus $\gamma$ does not have zero girth.  Lastly, the curve $\gamma$ cannot have girth one by Lemma~\ref{genus two}.  So we are done with this implication.

\textbf{GC $\Rightarrow$~\ref{ag girth conjecture}:}
Now assume the Girth Conjecture for $H\cup H^\prime$.  It suffices to show that either $H\cup H^\prime$ is a genus two splitting of a genus two spherical manifold, or the algebraic and geometric girths of $H\cup H^\prime$ are equal.

If the genus of $H$ is one, then $H\cup H^\prime$ is a standard Heegaard splitting of a Lens space.  It is straightforward in this case that the algebraic and geometric girths are equal.  So assume that the genus of $H$ is at least two.

We have $\pi_1(M)=\pi_1(H)/\rr$, where $\rr$ is the group of relators.  If $H\cup H^\prime$ is a reducible Heegaard splitting, then the geometric girth of 
$H\cup H^\prime$ is zero.  In this case, the algebraic and geometric 
girths must be equal.

Assume then that $H\cup H^\prime$ is an irreducible Heegaard splitting. By~\cite{Ha}, the 3-manifold $M$ is irreducible, thus $\rr$ is nontrivial.  Let $\gamma$ be a closed curve in $\bd H$ that realizes the algebraic girth of the splitting $H\cup H^\prime$.  Then $\gamma$ is a nontrivial element of $\rr$, and $\gamma$ bounds a disc in $H^\prime$.  By Corollary~\ref{easy inequality}, the algebraic girth of $H\cup H^{\prime}$ is less than or equal to the geometric girth.    To show the algebraic and geometric girths are equal, it suffices to show the geometric girth is less than or equal to $\girth(\gamma)$.

Assume temporarily that $\gamma$ is primitive in $\pi_1(H)$.  Apply the Girth
Conjecture to $\gamma$.  Then there exists a simple closed curve $\gamma^\prime$ in $\bd H$ that bounds a disc in $H^\prime$ and satisfies $\girth(\gamma^{\prime})\leq \girth(\gamma)$.  Since the geometric girth of $H\cup H^\prime$ is less than or equal to $\girth(\gamma^\prime)$, we know the geometric girth is less than or equal to $\girth(\gamma)$.  So we are done in this case.

Assume from now on that $\gamma$ is not primitive in $\pi_1(H)$.  Then 
$\gamma$ is homotopic to $\alpha^n$ for some primitive $\alpha$ in $\pi_1(H)$.  By Lemma~\ref{power},
$\girth(\alpha)\leq\girth(\gamma)$.  If $\alpha$ is also in $\rr$, we 
can apply the Girth Conjecture to $\alpha$ as in the previous 
paragraph to be done.  So assume that $\alpha$ does not lie in $\rr$.
Then $\pi_1(M)$ has torsion.  Since $M$ is irreducible, we must 
have that $\pi_1(M)$ is finite~\cite[Chapter 9]{He1}.

We now treat the cases of $\genus(H)=2$ and $\genus(H)>2$, separately.
Assume first that $\genus(H)=2$.  By the work of Thurston, genus two manifolds are known to satisfy the Geometrization Conjecture~\cite{Thu1}. Since $\pi_{1}(M)$ is finite, the manifold $M$ is a spherical 3-manifold. Since $H\cup H^\prime$ is irreducible, the splitting must be a genus two splitting of a genus two spherical manifold.

Assume finally that $\genus(H)\geq 3$.  Represent $\pi_1(H)$ as a free 
group $\langle a, b, c, \cdots\rangle$ on finitely many generators
by choosing a disc system for $H$ and letting each disc correspond to
a generator.  Since $\pi_1(M)=\pi_1(H)/\rr$ is finite, there exist 
$m,n\geq 1$ so that both $a^m$ and $b^n$ are trivial in $\pi_1(M)$.  
Then $a^mb^n$ is primitive element in $\pi_1(H)$ and trivial in $\pi_{1}(M)$.  Since $a^mb^n$ can be represented by a curve in $H$ disjoint from 
the disc corresponding to $c$, the element $a^m b^n$ has girth zero.

Represent $a^mb^n$ by a curve in $\bd H$ that contracts in $H^\prime$.  Applying the Girth Conjecture to this curve, there exists a simple closed curve in $\bd H$ of girth zero in $H$ that bounds a disc in $H^\prime$.  In other words, the splitting $H\cup H^\prime$ has geometric girth zero.  By Corollary~\ref{easy inequality}, $H\cup H^\prime$ must also have algebraic girth zero.  The algebraic and geometric girths of $H\cup H^{\prime}$ are equal, so we are done.
\end{proof}